% 第 22 讲 阿廷与布劳尔互反律（下）（源：src/lecture22.txt；已删去抽出时残留的页眉「22. ARTIN AND BRAUER RECIPROCITY, PART II 97/99」）
\section*{第 22 讲\quad 阿廷与布劳尔互反律（下）}

本讲的目标是在模去第二不等式的条件下证明阿廷互反律与布劳尔互反律。回忆阿廷互反律的陈述：对数域 $F$ 与素位 $v$，lcft 给出映射
\[
\theta_v:F_v^\times\to\Gal^{\mathrm{ab}}(F),
\]
而我们断言诱导映射
\[
\theta:A_F^\times\to\Gal^{\mathrm{ab}}(F)
\]
限制到 $F^\times$ 上平凡。这是二次互反律的巨大推广：若 $F=\Q$，则对任意二次扩张 $\Q(\sqrt{\ell})$，有 $\theta_v(\cdot)=(\cdot,\ell)_v$，即 $F_v$ 上的希尔伯特符号，其中 $v$ 是素数或 $\infty$。
% [原文作 "for any quadratic extension Q(`)"，抽出时丢了根号；按上下文还原为 Q(√ℓ)，见报告 §3]
然而，这里我们断言这对整个交换化的 Galois 群成立，而不只是对这个同构于 $\Z/2\Z$ 的特殊商成立。

\textbf{证明（阿廷互反律）.} 首先注意，对任意阿贝尔扩张 $E/F$，我们得到映射 $\theta:A_F^\times\to\Gal(E/F)$；显然只需在每个这样的商中证明消没。

\textbf{情形 1.} 设 $F:=\Q$，$E:=\Q(\zeta_{\ell^r})$，其中 $\ell$ 是素数。必须证明对每个素数 $p$ 有 $\theta(p)=1$，且 $\theta(-1)=1$，因为这些元素生成 $\Q^\times$ 而 $\theta$ 是群同态。我们借助 Dwork 定理作显式计算。先设 $p\ne\ell$。则
\[
\begin{cases}
\theta_q(p)=1,\\
\theta_\infty(p)=1\\
\theta_p(p)=p,\\
\theta_\ell(p)=p^{-1},
\end{cases}
\qquad\text{其中 } q\ne p,\ell.
\]
对第一个等式，注意只要 $q\ne\ell$，分圆扩张 $E/F$ 在 $q$ 处非分歧，而局部阿廷映射在 $q$ 处非分歧的扩张上杀掉 $(\Z/q\Z)^\times$ 的每个元素（见断言 21.5 的证明）。对第二个等式，注意每个 $\theta_v$ 都经过 $v$ 处的分解群分解，而此情形下分解群是 $\Z/2\Z$，且不难显式看出 $\theta_\infty:\R^\times\to\Z/2\Z$ 对应符号函数。对第三个等式，注意 $E/F$ 在 $p$ 处非分歧，因为 $p\ne\ell$。由 Dwork 定理，素元 $p$ 映到 $\Gal(E/F)=(\Z/\ell^r\Z)^\times$ 中的 Frobenius 元素，而它就是 $p$。最后，$E/F$ 在 $\ell$ 处完全分歧，且 $p\in(\Z/\ell\Z)^\times$，故 Dwork 定理给出 $\theta_\ell(p)=p^{-1}$。

现在设 $p=\ell$。则
\[
\begin{cases}
\theta_q(\ell)=1,\\
\theta_\infty(\ell)=1\\
\theta_\ell(\ell)=1,
\end{cases}
\qquad\text{其中 } q\ne\ell.
\]
最后一个等式由 Dwork 定理得到，因为 $\ell$ 在扩张的完全分歧因子上平凡作用（而在非分歧因子上作为其 Frobenius 元素作用）。

最后检验 $\theta(-1)$：
\[
\begin{cases}
\theta_q(-1)=1,\\
\theta_\infty(-1)=-1\\
\theta_\ell(-1)=(-1)^{-1},
\end{cases}
\qquad\text{其中 } q\ne\infty,\ell.
\]
由于 $-1\in(\Z/\ell\Z)^\times$，最后一种情形可如前应用 Dwork 定理。

\textbf{情形 2.} 仍设 $F:=\Q$，$E:=\Q(\zeta_n)$，$n$ 为任意整数。则
\[
n=\prod_{i=1}^m p_i^{r_i}
\]
其中 $p_i$ 为素数，因此
\[
\Q(\zeta_n)=\Q(\zeta_{p_1^{r_1}})\cdots\Q(\zeta_{p_m^{r_m}})
\]
是其各素数幂因子对应的复合，从而
\[
\Gal(\Q(\zeta_n)/\Q)=\prod_{i=1}^m\Gal(\Q(\zeta_{p_i^{r_i}})/\Q).
\]
于是阿廷映射
\[
\theta:A_\Q^\times\to\Gal(\Q(\zeta_n)/\Q)
\]
由 lcft 是各 $\Q(\zeta_{p_i^{r_i}})/\Q$ 的阿廷映射之积给出的，故只需注意下述一般性断言：线性不交扩张的阿廷互反律蕴含它们复合的阿廷互反律。也就是说，$\Q^\times$ 被杀死，因为它的每个坐标都被前一情形杀死。

\textbf{情形 3.} 设 $F$ 是一般数域，$E:=F(\zeta_n)$，$n$ 为某整数。由 lcft（在局部域乘法群的层面上），我们有下述交换图：
\[
\begin{CD}
F^\times @>>> A_F^\times @>{\theta}>> \Gal(F(\zeta_n)/F)\\
@V{N}VV @V{N}VV @VVV\\
\Q^\times @>>> A_\Q^\times @>{\theta}>> \Gal(\Q(\zeta_n)/\Q),
\end{CD}
\]
% [原文作 "at the level of multiplicative groups os local fields"，"os" 疑为 "of" 之误，见报告 §3]
其中 $\theta$ 表示阿廷映射。由于最右的映射是单射，只需证明 $N(F^\times)\subseteq\Q^\times$ 在 $\Gal(\Q(\zeta_n)/\Q)$ 中消没，而这恰是前一情形。

\textbf{情形 4.} 设 $E/F$ 是数域的分圆扩张，即对某个 $n$ 有 $E\subseteq F(\zeta_n)$。则由 lcft 有交换图
\[
\begin{array}{ccc}
A_F & \xrightarrow{\ \theta\ } & \Gal(F(\zeta_n)/F)\\[5pt]
& \stackrel{\theta}{\searrow} & \big\downarrow\\[5pt]
& & \Gal(E/F),
\end{array}
\]
而既然 $\theta$ 在上方的 Galois 群中杀掉 $F^\times$，它在下方那个中也如此。

我们已经证明阿廷互反律对所有数域的分圆扩张成立；对一般情形，我们将使用布劳尔互反律，现在转向它。\hfill$\square$

作为附注，Sam 大概是世上唯一称之为「布劳尔互反律」的人；通常它被称为「从 gcft 计算布劳尔群」之类。

设 $E/F$ 是整体域的 $G$-Galois 扩张。我们有 $G$-模的短正合列
\[
1\to E^\times\to A_E^\times\to C_E\to 1,
\]
它给出复合
\[
\begin{aligned}
\underbrace{\Br(F/E)}_{H^2(G,E^\times)}&\longrightarrow\underbrace{\bigoplus_v\Br(F_v/E_w)}_{H^2(G,A_E^\times)}\hookrightarrow\bigoplus_v\Br(F_v)\\
&\xrightarrow{\ (x_v)_{v\in M_F}\mapsto\sum_{v\in M_F\setminus M_F^\infty}x_v\ }\Q/\Z,
\end{aligned}
\]
其中回忆 $\Br(F/E)$ 表示 $F$ 上那些与 $E$ 张量后成为矩阵代数的可除代数构成的群；$w$ 是 $v$ 上方某个选定的素位；再回忆
\[
\Br(F_v)=
\begin{cases}
\Q/\Z & \text{若 } v \text{ 有限},\\
\tfrac12\Z/\Z & \text{若 } v \text{ 实},\\
0 & \text{若 } v \text{ 复};
\end{cases}
\]
而最右的映射被称为「不变量」映射，因为它是「局部不变量」之和，这些局部不变量分类局部可除代数。此复合对应于把中心单代数与各局部素位上的可除代数作张量、在每一这样的素位取不变量、再把它们相加。直和（自动地）告诉我们：除有限多个素位外，我们在域上得到的是矩阵代数。布劳尔互反律断言此复合为零。于是例如，不存在 $\Q$（或任意数域）上的中心单代数，使得它在除一个素位以外的所有素位上都是（分裂的）矩阵代数。注意当此复合作用于（由两个有理数给出的）哈密顿代数时，它不过是记录了希尔伯特互反律。

\textbf{断言 22.1.} 设 $E/F$ 是整体域的循环扩张。则阿廷互反律等价于布劳尔互反律。

\textbf{证明.} 取定一个生成元，有 $\Gal(E/F)=G\cong\Z/n\Z$。由于 $G$ 循环，它等于自身的交换化，且 Tate 上同调是 $2$-周期的，于是有下述交换图：
\[
\begin{CD}
\widehat H^0(G,E^\times)=F^\times/N(E^\times) @>>> \widehat H^0(G,A_E^\times)=A_F^\times/N(A_E^\times) @>{\theta}>> G=\tfrac1n\Z/\Z\\
@| @| @VVV\\
\Br(F/E)=\widehat H^2(G,E^\times) @>>> \widehat H^2(G,A_E^\times) @>{\iota}>> \Q/\Z,
\end{CD}
\]
其中 $\iota$ 表示不变量映射。现在，左边的方块平凡地交换，而右边的方块由循环扩张的 lcft 交换。断言随即由简单的图追踪得到。\hfill$\square$

\textbf{断言 22.2.} 对任意整体域 $F$ 与每个 $\beta\in\Br(F)$，存在循环分圆扩张 $E/F$ 使 $\beta\in\Br(F/E)$，即 $\beta$ 也落在相对布劳尔群中。

暂且假设此断言。则我们可以一般地推出布劳尔互反律：由于扩张 $E/F$ 循环，我们知道阿廷互反律与布劳尔互反律等价，而且我们还知道阿廷互反律成立，因为该扩张是分圆的。
% [原文作 "we know that Brauer reciprocity is equivalent to Brauer reciprocity"，前一处的 Brauer 应为 Artin，见报告 §3]

此断言对局部域容易验证：给定局部域 $K$ 上的可除代数，若其次数的平方根整除 $[L:K]$，则它在域扩张 $L/K$ 上分裂。事实上，回忆若 $\beta\subseteq\frac1n\Z/\Z\subseteq\Br(K)$，即 $\beta$ 是 $n^2$ 次可除代数，则 $\beta\in\Br(K/L)$ 当且仅当 $n\mid[L:K]$，因为
\[
\Q/\Z=\Br(K)\xrightarrow{\ \times[L:K]\ }\Br(L)=\Q/\Z.
\]

现在回忆下述定理，我们将在下一讲证明：

\textbf{定理 22.3.} 对整体域的一切扩张 $E/F$，有 $H^1(G,C_E)=0$。

这算是希尔伯特定理 90 的类比，而证明它是 gcft 中最困难的部分；我们暂且假设它。短正合列
\[
1\to E^\times\to A_E^\times\to C_E\to 1
\]
于是给出正合列
\[
\underbrace{H^1(A_E^\times)}_{0}\to\underbrace{H^1(C_E)}_{0}\to H^2(E)=\Br(F/E)\to H^2(A_E^\times)=\bigoplus_v\Br(F_v/E_w),
\]
% [原文此二项下有 underbrace，抽出时标注丢失；按紧随其后的「the vanishing is by the previous theorem and Hilbert's Theorem 90」补作 0，见报告 §2]
其中消没由前一定理与希尔伯特定理 90 给出（因为阿代尔的上同调就是各局部上同调的直和）。取正向极限，得到下述推论：

\textbf{推论 22.4.} 对任意整体域 $F$，有典范的单射
\[
\Br(F)\hookrightarrow\bigoplus_{v\in M_F}\Br(F_v).
\]

换言之，数域 $F$ 上的中心单代数是矩阵代数，当且仅当它在 $F$ 的每个完备化处都是矩阵代数。这绝不是显然的陈述！

要点是：为证明断言 22.2，只需证明下述断言：

\textbf{断言 22.5.} 给定整体域 $F$ 的有限素位集 $S$，以及对所有 $v\in S$ 的正整数 $m_v$（要求：$v$ 复时 $m_v=1$；$v$ 实时 $m_v=1,2$；$v$ 有限时 $m_v$ 任意），则存在循环分圆扩张 $E/F$，使得 $E_w/F_v$ 的次数被 $m_v$ 整除（对任意选定的 $w\mid v$）。
% [原文作 "such that m_v = 1 is v is complex"，第二个 "is" 疑为 "if" 之误，见报告 §3]

\textbf{证明（断言 22.5 $\Rightarrow$ 断言 22.2）.} 由推论 22.4，$F$ 上的中心单代数分裂，当且仅当它在每个局部域 $F_v$ 上分裂，而这又当且仅当其次数的平方根整除某个扩张 $E/F$ 与素位 $w\mid v$ 的 $[E_w:F_v]$。于是，对每个素位 $v$ 取 $m_v$ 为其在 $F_v$ 上次数的平方根（等价地，取其在 $v$ 处局部不变量的分母），当 $v$ 实时它必为 $1$ 或 $2$、$v$ 复时为 $1$，此断言蕴含我们的中心单代数在它所提供的扩张 $E$ 上分裂。此外，如前所述，这样的中心单代数在除有限多个素位 $v$ 外的所有 $v$ 上都在 $F_v$ 上分裂，这意味着我们可以取 $S$ 为仅由那些它不分裂的素位组成的集合。\hfill$\square$

在开始证明之前，让我们花点时间注意：确实存在不是循环的分圆扩张！事实上，$\Gal(\Q(\zeta_n)/\Q)=(\Z/n\Z)^\times$，其中 $\zeta_n$ 是本原 $n$ 次单位根，因此例如当 $n=15$ 时，
\[
(\Z/15\Z)^\times\cong(\Z/3\Z)^\times\times(\Z/5\Z)^\times\cong\Z/2\Z\times\Z/4\Z
\]
不是循环的！

\textbf{证明（断言 22.5）.} 首先注意可设 $F=\Q$：我们可把每个局部因子 $m_v$ 换为 $[F_v:\Q_p]\cdot m_v$（其中 $p$ 是 $v$ 下方的 $\Q$ 的素位），在 $\Q$ 上应用本断言，然后把 $E$ 换为复合 $E\cdot F$，它也是 $F$ 的循环分圆扩张（因为循环群的子群循环）。

\textbf{情形 1.} 设对每个局部素位 $v\in S$，有 $m_v=p^{r_v}$，其中 $p$ 是（与 $v$ 无关的）奇素数。对任意 $r$，有分圆扩张 $\Q(\zeta_{p^r})/\Q$，Galois 群为 $(\Z/p^r\Z)^\times$。注意 $\#(\Z/p^r\Z)^\times=(p-1)p^{r-1}$，故可令 $E_r$ 为 $\Q(\zeta_{p^r})/\Q$ 的循环子扩张，$\Gal(E_r/\Q)\cong\Z/p^{r-1}\Z$。则对每个 $v\in S$ 与任意选定的 $w\mid v$，我们断言 $[E_{r,w}:\Q_v]\to\infty$（$r\to\infty$）。事实上，除 $\C$ 外的任何局部域只含有限多个单位根，故这些扩张的次数必递增。现在，
\[
[E_{r,w}:\Q_v]=[\Q_v(\zeta_{p^r}):\Q_v]\mid(p-1)p^\infty,
\]
% [原文抽出作 "[Q_v(ζ_{p^r}/Q_v] | (p−1)p^∞"，缺右方括号与冒号；按图像排作 [Q_v(ζ_{p^r}):Q_v]，见报告 §3]
故此 $p$ 幂因子发散；由于 $S$ 有限，本情形得证。请注意这些扩张是全复的（即每个无限素位都是复的），故不必担心那些 $m_v=2$ 的实素位 $v$。

\textbf{情形 2.} 现设对每个 $v\in S$ 有 $m_v=2^{r_v}$。此情形类似，只差一点：$(\Z/2^r\Z)^\times\cong\Z/2^{r-2}\Z\times\Z/2\Z$，此同构可通过验证下述容易的恒等式得到：
\[
\Z/2^{r-2}\Z\cong\{x\in(\Z/2^r\Z)^\times:x\equiv1\bmod4\}
\]
（于是可通过计算其 $2$-挠部分、再验证它只含 $2$ 个元素等等，来证明它是循环的），并注意 $\Z/2\Z\cong\{1,2^{r-1}-1\}$。于是可令 $\Q(\zeta_{2^r})/E_r/\Q$ 为分圆扩张，其次数为 $2^{r-2}$，Galois 群为 $(\Z/2^{r-2}\Z)^\times$，实现为 $(\Z/2^r\Z)^\times$ 的商。
% [原文作 "a cyclotomic degree-2^{r-2} extension whose Galois group is (Z/2^{r-2}Z)^×"，(Z/2^{r-2}Z)^× 的阶为 2^{r-3}，与次数 2^{r-2} 不符；疑笔误，照译并见报告 §3]

我们断言 $E_r/\Q$ 是全复扩张。由于它是阿贝尔扩张，只需证明复共轭（对应 $-1\in(\Z/2^r\Z)^\times$）在其上非平凡地作用。对充分大的 $r$，$-1\not\equiv1\bmod4$，故由上述显式同构，它在 $(\Z/2^{r-2}\Z)^\times$ 中的投影为 $1-2^{r-1}\ne1$（注意 $2^{r-1}-1\equiv3\bmod4$ 是任一因子中唯一具有此性质的非平凡元素）。于是，由与前一情形类似的论证，扩张 $E_r/\Q$ 即已足够。

\textbf{情形 3.} 最后，对一般情形，取前两种情形所构造的、对应于 $m_v$ 的各素因子的扩张的复合。此外不会发生任何「干扰」（以致复合不再循环），因为每个扩张的 Galois 群都有互不相同的素数幂阶。\hfill$\square$

于是，假设我们关于希尔伯特定理 90 的类比成立，整体域 $F$ 的布劳尔群的每个元素都被某个循环分圆扩张分裂，这蕴含对 $\Br(F)$ 所有元素的布劳尔互反律。这又蕴含循环扩张的阿廷互反律，因为对循环扩张这两条互反律等价。最后，由于任意阿贝尔群都是循环群之积，每个阿贝尔扩张 $E/F$ 都是循环扩张的复合，从而一般地蕴含阿廷互反律（我们已经看到，一组线性不交扩张的阿廷互反律蕴含它们复合的阿廷互反律）。
